\begin{lemma}[SAP forces a common infinite invariant monolith]
\label{lem:rh-sap-invariant-monolith}
Let $M$ be an INFINITE transitive homogeneous structure
with SAP, in a finite COMPLETE BINARY original language.
No symmetry, CCO, distributivity or triangle constraints
are needed. For any TWO nonidentity original invariant
equivalences $E,F>\Id$, their actual intersection
$E\cap F$ has INFINITE classes.

Consequently the finite FULL invariant-equivalence lattice
has a UNIQUE least nonidentity element $L>\Id$, whose
classes are infinite. The greatest finite-class original
equivalence is $F_{\mathrm{fin}}=\Id$; in the notation of
\cref{cor:rh-acl-bounded-character} there is exactly
ONE minimal infinite direction ($r=1$), including
$L=\Top$ when the source is primitive.
Moreover $\operatorname{acl}_M(A)=A$ for EVERY finite
ORIGINAL REAL $A$. These SAP conclusions do NOT
hold for AP-only sources such as $M_r=\mathbb N^r$.
\end{lemma}
\begin{proof}
Fix an actual real $x$. Since $E,F$ are nonidentity
and the original sort is TRANSITIVE,
choose $y\ne x$ with $xEy$ and
$z\ne x$ with $xFz$.
The base $A=\{y,z\}$ is a COMPLETE induced original
finite substructure (possibly one point if $y=z$),
and $B=A\cup\{x\}$ is a complete extension with
$x\notin A$.  For ANY positive integer $n$,
iteratively STRONGLY amalgamate $n$ copies of $B$
over the same labelled copy of $A$.
There are pairwise distinct new real points
$x_1,\ldots,x_n$, each satisfying
$x_iEy$ and $x_iFz$. Consequently
\[
                 x_i\ (E\cap F)\ x_j
                  \qquad(1\le i,j\le n).
\]
The finite amalgam embeds in $M$ by the age
property; as $n$ was arbitrary, some
$E\cap F$-class is infinite, and transitivity
of the original sort makes ALL its classes
infinite. In particular $E\cap F>\Id$.

There are only FINITELY many original invariant
equivalences, since each is a union of
finitely many COMPLETE original pair orbitals.
Iterating the preceding result shows that the
intersection of ALL nonidentity invariant
equivalences is still nonidentity.
Call it $L$. Then $L$ is the UNIQUE least
nonidentity invariant equivalence and has
infinite classes. Every nonidentity invariant
equivalence has infinite classes (take $E=F$),
so the greatest finite-class original equivalence
is exactly $\Id$. Thus the sole minimal infinite
equivalence is $L$ and $r=1$.

Finally fix a finite original real base $A$
and a real point $u\notin A$.
By repeatedly strongly amalgamating
the COMPLETE extension $A\subseteq A\cup\{u\}$
over the complete base $A$, then embedding
the result into $M$ and aligning its base by
homogeneity, we obtain arbitrarily many
realizations of the COMPLETE old one-point
type of $u$ over $A$. Homogeneity makes them
one infinite $\Aut(M/A)$-orbit.
Thus no off-base real is algebraic, and
$\operatorname{acl}_M(A)=A$.
\end{proof}

\begin{corollary}[Infinite joint real fibre of all proper MI quotient coordinates]
\label{cor:rh-sap-mi-joint-real-fibre}
Let $M$ be an INFINITE transitive homogeneous
finite COMPLETE BINARY original structure
with SAP, with no symmetry or CCO assumption.
Let $L>\Id$ be its unique least nonidentity
ORIGINAL invariant equivalence from
\cref{lem:rh-sap-invariant-monolith}.
Let $\mathcal I^\times$ be the set of
ALL proper NONIDENTITY meet-irreducible
equivalences in the FULL original invariant
lattice. (The original REAL identity is
EXCLUDED, even if it is meet-irreducible.)
With the empty meet interpreted as $\Top$,
\[
           \boxed{\quad
                L=\bigcap_{J\in\mathcal I^\times}J.
                 \quad}
 \tag{MK}\label{eq:rh-sap-mi-meet}
\]
Thus the ACTUAL map on original real points
\[
          \Phi(x)=([x]_J)_{J\in\mathcal I^\times}
\]
has infinite real fibres:
\[
        \Phi^{-1}(\Phi(x))=[x]_L,\qquad
        \left|\Aut(M/\Phi(x))\cdot x\right|
                    =|[x]_L|=\infty.
\]
If $\mathcal I^\times=\varnothing$
(the primitive case), then $L=\Top$,
the empty tuple is fixed, and the sole
fibre is all of $M$.  The statement
concerns an ACTUALLY REALISED full tuple
of MI quotient values, not an arbitrary
tuple of classes. It does not assert
SAP for a hidden language naming
additional non-MI exact quotients.
\end{corollary}
\begin{proof}
In any FINITE lattice, each element
$X$ is the meet of all proper
meet-irreducibles $J\ge X$.
Indeed, for any $Y\nleq X$, choose
a maximal $J\ge X$ with $Y\nleq J$.
If $J=J_1\cap J_2$ with both
$J_i>J$, maximality would imply
$Y\le J_1,J_2$, hence $Y\le J$,
a contradiction. Thus $J$ is
meet-irreducible and separates
$Y$ from $X$, proving the finite
lattice identity without any
distributivity hypothesis.

Apply it to $X=L$.  Since $L$
is the UNIQUE least nonidentity
invariant equivalence, EVERY
nonidentity invariant equivalence
contains $L$. Consequently all
proper meet-irreducibles above $L$
are exactly $\mathcal I^\times$,
and their meet is $L$.
The finite lattice includes the
primitive empty-list convention.

Now $y$ has the SAME entire MI tuple
as $x$ iff $xJy$ for every
$J\in\mathcal I^\times$, equivalently
$xLy$. Thus its fibre is $[x]_L$,
which is infinite by
\cref{lem:rh-sap-invariant-monolith}.
The pointwise stabiliser of this
tuple is the setwise stabiliser of
$[x]_L$: an automorphism fixes all
$J$-classes of $x$ iff its image
of $x$ lies in their intersection.
Original transitivity then makes
this stabiliser TRANSITIVE on
$[x]_L$, yielding the asserted
infinite orbit. No REAL equality
is represented by a positive
diversity factor in this proof.
\end{proof}

\begin{corollary}[Infinite Helly intersections of proper original exact classes under SAP]
\label{cor:rh-sap-infinite-exact-helly}
Let $M$ be an infinite transitive homogeneous
finite complete SYMMETRIC BINARY original
structure with SAP and semantic
HOMOMORPHIC CCO. Let
$E_1,\ldots,E_m>\Id$ be any nonidentity
FULL RAW exact original invariant
equivalences on the SAME original
real sort, and choose actual
classes $C_i=[x_i]_{E_i}$.
If all $C_i$ meet pairwise, then
\[
    \boxed{\quad\bigcap_{i=1}^m C_i
                  \ \text{is INFINITE}.\quad}
\]
More precisely, with
$L=\bigcap_{i=1}^m E_i$,
the intersection is ONE entire
$L$-class; $L>\Id$ and every
$L$-class is infinite.
If in addition the source is TERMINAL
(full invariant-equivalence NR2), the
result holds for ALL nonidentity
original invariant equivalences:
no separate exactness assumption is
needed by
\cref{lem:typed-terminal-all-equivalences-exact}.
For $m=0$ the intersection is the
whole infinite original sort.
For an identity kernel the infinite-size
conclusion is NOT asserted.

In particular, in the terminal setting,
a FULL tuple of proper nonidentity
meet-irreducible original quotient
classes is realised by an original
REAL point if and only if the classes
are PAIRWISE compatible. Whenever
realised, their entire common
real fibre is precisely one
INFINITE class of the unique original
SAP monolith $L$ by
\cref{cor:rh-sap-mi-joint-real-fibre}.
This is an original real-incidence
statement and does NOT imply SAP
after naming arbitrary NON-MI
exact meet quotient sorts.
\end{corollary}
\begin{proof}
For $m\ge1$, the actual classes
$C_i=B_{E_i}(x_i)$ are FULL RAW exact
balls, and pairwise intersection is
the exact hypothesis of the old
exact-ball $2$-Helly theorem
\cref{thm:full-typed-exact-balls}(iii)
on the single original sort.
Hence choose $x\in\bigcap_i C_i$.
Since $x\in C_i$, one has
$C_i=[x]_{E_i}$ for EVERY $i$.
Therefore
\[
               \bigcap_i C_i
                    =[x]_{\bigcap_iE_i}.
\]
By \cref{lem:rh-sap-invariant-monolith},
the intersection of FINITELY MANY
nonidentity invariant equivalences
remains nonidentity with infinite
classes. Thus the displayed set
is infinite, as claimed.
When full NR2 is also assumed,
\cref{lem:typed-terminal-all-equivalences-exact}
shows every original invariant
equivalence is FULL RAW exact;
the same argument then applies
to all nonidentity invariant
equivalence classes.
\end{proof}

\begin{example}[Two-coordinate real algebraic corners]\label{ex:two-coordinate}
The absolute real algebraic closure in the hull theorem need not be
$F[A]$, even when the finite-class packet bottom $F$ is regular and
thin.  Let $G=\mathbb F_2$ and take
$M_*=\mathbb N\times\mathbb N\times G$.  For two distinct triples
$(r,c,g)$ and $(r',c',g')$, name the fine colour by the three data
$(\mathbf 1_{r=r'},\mathbf 1_{c=c'},g+g')$.  The only missing triple
$(1,1,0)$ denotes equality, so the diversity palette has seven
colours.  Every named colour is symmetric.

This is a transitive homogeneous, triangle-defined AP source with
semantic homomorphic CCO.  Indeed finite partial isomorphisms extend
by independently permuting the two infinite coordinates and by a
single translation of $G$.  In a complete diagram, legality of all
triangles makes each coordinate-equality relation transitive and makes
the prescribed $G$-differences a coboundary; the absent equality colour
prevents two distinct vertices from receiving the same triple.  For a
partial diagram, each specified edge imposes equality or inequality in
each infinite coordinate and one $G$-difference.  Inconsistency of
an equality/inequality system has an inequality edge joined by a path
of equality edges; inconsistency of the $G$-differences has a cycle
with nonzero total difference.  Each obstruction contains a specified
simple cycle without a homomorphism.  Conversely, in the absence of
such obstructions the equality components can be assigned distinct
coordinate values and the $G$-differences integrated.  This proves
semantic homomorphic CCO, not merely ordinary AP.

Here $F$ consists of the two-point packets
$\{(r,c,0),(r,c,1)\}$.  For
$a=(0,1,0)$, $b=(1,0,0)$ and $x=(0,0,0)$, the two colours
$\col(x,a)=(1,0,0)$ and $\col(x,b)=(0,1,0)$ determine $x$
uniquely.  Hence
\[
 x\in\operatorname{dcl}_M(\{a,b\})\setminus F[\{a,b\}].
\]
The two infinite exact equivalences generated by these colours are,
respectively, equality of the first coordinate and $G$-coordinate,
and equality of the second coordinate and $G$-coordinate.  Their
intersection is equality; the finite intersection in the hull theorem
is exactly $\{x\}$.  Thus the full canonical-hull argument is genuinely
stronger than closure under finite packets.  This example does not
contradict the \emph{relative} owner-closure theorem below, where
additional owner imaginaries are named.
\end{example}


\begin{remark}[The thin phase packet need not be exact]
\label{rem:phase-thin-not-exact}
In the same Cartesian phase model, the two-element packet
$F[(r,c,g)]=\{(r,c,0),(r,c,1)\}$ is a raw invariant
finite-class equivalence but is \emph{not} exact.  Indeed each
raw atom specifies a fixed phase difference in $\mathbb F_2$.
Consequently every nonempty atomic word profile has a single
fixed endpoint phase difference, equal to the sum of the
letters' phase differences.  An intersection of word profiles
has the same property whenever it is nonempty.  The only exact
support obtained from the empty defining family is $\Top$.
Since the proper equivalence $F$ contains both equality
(phase difference $0$) and the thin flip orbital (phase
difference $1$), it is not exact.  The same parity argument
excludes the invariant equivalences ``same row irrespective of
phase'' and ``same column irrespective of phase'' from the
proper exact supports.  This is a concrete reason for retaining
\emph{exactness} in the quotient theorems; arbitrary raw invariant
quotients are not interchangeable with exact quotients.
\end{remark}

\begin{remark}[Scope of the CCO substitution]
The proof applies to the \emph{actual real} structure $M$ and its realised
one-point fibres over finite real sets. It does not prove the stronger
AP-only version of the hull theorem or assert CCO inheritance to arbitrary
nonexact imaginary sorts. Corollary~\ref{cor:hasse}, rather than the
finite-group classification used in the original AP-only proof, is the
entire finite-cover input here.
\end{remark}
